\documentclass[11pt]{article}
\usepackage{amsmath,amssymb,hyperref}
\title{Independent Root Memories with Admissible Squared Cauchy Corrections}
\author{Joint Cell Head Compression}
\date{November 24, 2026}
\begin{document}
\maketitle
\paragraph{Status.} The identities and two-head implementation below are
kernel checked. They do not prove a new asymptotic bound on worst-case head
complexity. The original model and all of its definitions are unchanged.

\section{The numerator gap and its repair}
For independent roots $u_1,\ldots,u_c$, partial fractions give
\[
 \frac1{\prod_l(p-u_l)}=\sum_l\frac{a_l(u)}{p-u_l},
 \qquad a_l(u)=\prod_{m\ne l}(u_l-u_m)^{-1}.
\]
The roots may vary independently from column to column. Multiplying this
identity by a row-dependent numerator $A_i$ does not directly give $c$
admissible heads: the coefficient $A_i a_l(u_j)$ is jointly dependent on
row and column, whereas a head numerator must be affine on the whole
Boolean cube.

Differentiation in $p$ gives the alternative kernel
\[
 \frac{P_u'(p)}{P_u(p)^2}
 =\sum_l\frac{a_l(u)}{(p-u_l)^2},\qquad P_u(p)=\prod_l(p-u_l).
\]
Each weighted term now has a two-head implementation. For arbitrary
$A_i,c_j$ and compact offsets, set
\[
 Q_t(i,j)=\frac{tA_i}{D_i-T_j-c_j/t}
          -\frac{tA_i}{D_i-T_j}
 =\frac{A_i c_j}{(D_i-T_j-c_j/t)(D_i-T_j)}.
\]
Thus $Q_t\to A_i c_j/(D_i-T_j)^2$. This is a different kernel from the
ordinary product inverse; it cannot be substituted into an existing
sign decoder without proving its signs again.

\section{Full-cube admissibility and exact axes}
Take $1-\delta<D_i<1$, $0<T_j<\delta$, and $(r+s)\delta<1$.
For all sufficiently large positive $t$, also
$0<T_j+c_j/t<\delta$. Each denominator on the entire Boolean cube is
\[
 1-\sum_i(1-D_i)x_i-\sum_j U_jy_j,
 \qquad U_j=T_j\text{ or }T_j+c_j/t.
\]
All slopes are strictly negative and the sum of their magnitudes is less
than one. Both numerators are affine row-only functions, with coefficients
$\pm tA_i$. When all row bits vanish, both numerators vanish. When all
column bits vanish, the two denominators coincide and the numerators cancel.
These are exact zeros on both entire Boolean axes at every finite stage,
not merely limits on one-hot inputs.

For $G$ unrelated terms, the budget is exactly $2G$. One sufficiently
large common $t$ preserves every specified strict signed margin after
addition to an arbitrary prior score on the finite grid. Applying this to
$c$ roots yields $2c$ heads per product-derivative correction, with roots
chosen independently in every column. No common label table is assumed.

\section{Checked declarations and remaining capacity problem}
The files are \texttt{JointCellSquaredCauchy.lean} and
\texttt{JointCellRootProduct.lean}, in namespace
\texttt{HeadComplexity.AxisZeroGluing}. The main declarations are
\texttt{exists\_squaredCauchyApprox\_data},
\texttt{squaredCauchy\_preserves\_signs},
\texttt{rootProduct\_partial\_fractions},
\texttt{rootProduct\_squared\_partial\_fractions}, and
\texttt{rootProduct\_preserves\_signs}.
Their transitive axioms are only \texttt{propext},
\texttt{Classical.choice}, and \texttt{Quot.sound}.

The row functions $A,D$ in the checked grid theorem are freely chosen
because the row is one-hot. On a larger product domain they must be affine
in its actual Boolean encoding, or a further compiler must be charged.
In particular, a bit-dependent tensor residue is not affine merely because
it is additive separately in the point coordinates. Replicating the bit
index across blocks changes both the ambient dimension and the support
problem. No universal vector-memory decoder with a growing net saving is
asserted. The required next result is a supported decoder with arbitrary
independent labels whose capacity exceeds its fully charged implementation
cost by an unbounded factor beyond the inherited logarithm.
\end{document}
